\section{A fixed current-syndrome decoder}\label{sec:decoder}

We construct a deterministic function \(D(s)\) with
\(\partial_2D(s)=s\) on every attainable syndrome \(s\). The same
definition is used on the primal and dual complexes, with separate
fixed orderings of their cells. No mark times enter this definition.

\subsection{Filling a small component}

Set \(\delta=L/100\). Fix representatives of torus vertices and a
lexicographic ordering of all torus plaquettes. Suppose every
edge-connected component of \(s\) has ambient torus diameter less
than \(\delta\). For a nonempty component \(\Gamma\), choose its
lexicographically least vertex \(v_\Gamma\). Lift every vertex by
its shortest coordinate displacement in \((-L/2,L/2]\) from this
anchor, and let \(K_\Gamma\) be the coordinate box between the lifted
coordinate extrema.

Each lifted vertex lies within \(\delta\) of the anchor. The box has
each side length at most \(2\delta\); its projection embeds for large
\(L\). Neighbouring vertices lift consistently: a coordinate
difference between their lifts is at most \(2\delta<L-1\), which
rules out the spurious displacement \(L-1\) across a unit torus edge.
All edges of \(\Gamma\) therefore lift to edges of this box.
Even degree at every vertex makes the component a cellular 1-cycle.

Here is an explicit reason that this cycle fills. A cubical box is
a product of integer intervals, allowing intervals of length zero.
Collapse the final layer in one nondegenerate coordinate interval
onto the preceding layer. On cellular chains let \(p\) be that
projection, and let \(\kappa_h\) send a horizontal cell in the final layer
to its prism across the collapsed unit interval; let \(\kappa_h\) vanish
on other cells. Directly on each cell,
\[
 \partial\kappa_h+\kappa_h\partial=\id+p
\]
over \(\F\): the two horizontal faces of a prism give the cell and
its projected copy, while the side-face terms cancel. For a cell
crossing the layer, its top horizontal face supplies that prism
in \(\kappa_h\partial\), giving the same identity. Iterating the collapses
in all coordinates contracts the box to a vertex. Thus every
positive-dimensional cycle in the box is a boundary, including
every 1-cycle and 2-cycle. This also covers degenerate boxes.

Choose the lexicographically least plaquette subset in \(K_\Gamma\)
whose boundary is \(\Gamma\), with ordering induced by its projection
to torus cells. A finite set of fillings is nonempty by the
contraction, so this choice is well defined. Project back to the
torus and add the component fillings modulo two. Call the resulting
chain \(D(s)\). Set \(D(0)=0\). If some component has diameter at
least \(\delta\), use instead the lexicographically least global
solution of \(\partial_2D=s\). Attainability means
\(s\in\operatorname{im}\partial_2\), so a solution exists in that
case as well.

The construction is fixed once the cell orderings are fixed. In
particular an unchanged small component has exactly the same filling
at every observation. This consistency, rather than optimal area,
is its essential property.

\subsection{Invariance across a physical face flip}

For a current error chain \(E\), define its filled homology class by
\begin{equation}\label{eq:filled-class}
 \mathfrak h(E)=[E+D(\partial_2E)]\in H_2(K_L;\F).
\end{equation}
The bracket is meaningful since the displayed chain is closed.

\begin{lemma}[Local preservation of filled homology]\label{lem:homology}
For all sufficiently large \(L\), let \(E'=E+p\) differ from \(E\)
by one plaquette. If every syndrome component of both
\(s=\partial_2E\) and \(s'=\partial_2E'\) has diameter less than
\(\delta\), then
\[
 \mathfrak h(E')=\mathfrak h(E).
\]
This holds for arbitrary merges, splits, creation or disappearance
of components at that flip, and on the dual complex as well.
\end{lemma}
\begin{proof}
The only modified syndrome edges lie in \(\partial p\).
Any component in either decomposition that does not meet a vertex
of \(p\) remains unchanged. Indeed its edges and their incident
connections are untouched; a changed connection would have an
endpoint on \(p\). Conversely, a component changed by deletion or
insertion has a path to a modified edge and hence meets a vertex
of \(p\). The fillings of all unchanged components cancel in
\(D(s)+D(s')\).

Lift \(p\) once to the integer lattice. Align the anchor lift of
each affected component using one of its vertices on this face.
Its anchor is within \(\delta\) of that shared vertex, and its
filling box is within \(\delta\) in each coordinate of the anchor.
Thus all affected fillings, together with \(p\), are supported in
one integer box within coordinate distance \(2\delta+2\) of a fixed
vertex of \(p\). For large \(L\) the width \(4\delta+4\) is less
than \(L-1\), so the common box embeds in the torus.

The chain
\[
 C=p+D(s)+D(s')
\]
has boundary \(\partial p+s+s'=0\). Its uncancelled support is in
the common box. The contraction above makes \(C\) the boundary of
a 3-chain there. Consequently adding \(C\) does not change the
homology class in \eqref{eq:filled-class}. Duality preserves all
these incidence and support properties, proving the last assertion.
\end{proof}

\begin{corollary}[Retention along the marked history]\label{cor:retention}
In either sector, started with \(E=0\), the probability that
\(\mathfrak h(E_s)\ne0\) at any \(s\le t\) is at most
\(C_0(1+t)e^{-\beta L}\), with constants as in
Proposition~\ref{prop:history}. An initial closed chain \(E_0\)
instead gives retention of its initial filled class with the same
bound.
\end{corollary}
\begin{proof}
On a finite torus there are finitely many clock rings up to each
finite \(t\), almost surely, even if that deterministic \(t\) is
exponentially large in \(L\). Order all effective face flips.
If the large-component event of
Proposition~\ref{prop:history} does not occur, each flip has only
small components on both sides. Lemma~\ref{lem:homology} propagates
the initial class through all of them. With \(E=0\) it starts at
zero. For an added closed initial chain, the syndrome and the
enabled correction predicates are identical under common marks;
the decoded class is shifted by its fixed initial class. The same
large-component event therefore controls a change in that class.
\end{proof}

The decoder need not infer an initial logical class from the
syndrome. An encoded start supplies that class. More generally, a
syndrome-only recovery can preserve a known encoded input through
this dynamics but cannot distinguish two initial errors differing
by a hidden logical membrane.
